% ---------------------------------------------------------------------------
\section{Introduction}
\label{sec:intro}

Interactive lenses are one of visualization's most general interaction
primitives: a movable region under which the base visualization is shown
differently~\cite{bier1993toolglass,tominski2014star,tominski2017survey}. On
maps, many lenses \emph{magnify}, \emph{filter} or \emph{re-render} what lies
beneath them. Others instead \emph{summarise} it: a probe pairs a
region with charts of its contents~\cite{butkiewicz2008probes}, a time lens
wraps a query circle in a ring of temporal aggregates~\cite{tominski2012stacking},
and a shipped map-lens product draws a bar chart of the data under each lens
around its rim~\cite{visquill}.

A summarising lens has to answer two questions that a magnifier does not. What
is summarised---which members, decomposed how, normalised against what? And
where does the summary go---in which order around the lens, on what curve, and
how is each mark tied back to the ground it describes? In the rings we have
seen, the second question is answered by habit: marks are laid out in category
(or time) order, so the most valuable channel on a ring---angular
position---carries a category or time order and nothing geographic. A bar at eleven
o'clock says nothing about the north-west.

This paper makes three moves. First, we refine Tominski et al.'s lens model
(selection $\Sel$, lens function $\lambda$, join
$\Jn$~\cite{tominski2017survey}) for summarising lenses into seven stages,
$\lambda=\Mrk\circ\Nrm\circ\Bin$ and $\Jn=\Asc\circ\Anc\circ\Plc$
(\cref{sec:model}), and use it to describe existing techniques
(\cref{sec:descriptive}) and to generate new ones (\cref{sec:generative}).
Second, we make angle mean \emph{bearing}: each aggregate is placed on the lens
boundary at the bearing of what it summarises, using necklace-map
placement~\cite{speckmann2010necklace} with minimal displacement; a leader is
drawn exactly where the gap between a mark and its data would mislead; and
because placement is solved in the curve's parameter, the ring unrolls into an
axis without re-solving. Third, we treat the number of lenses as a parameter,
so that one lens, a handful of small multiples and a gridded glyphmap are the
same object, and we are explicit about how much of that continuum is
geographically weighted statistics under another name (\cref{sec:discussion}).

\paragraph{Contributions.}
\begin{itemize}[leftmargin=*, itemsep=1pt]
  \item A seven-stage model of summarising lenses that refines the
    selection--function--join model, with two values the existing lens design
    spaces lack: an effect scope \emph{on the boundary}, and thematic marks
    under a spatial layout (\cref{sec:model}).
  \item Bearing-faithful placement of aggregates on a lens boundary, with a
    residual leader and curvature as a free axis (\cref{sec:placement,sec:association}).
  \item Within-unit structure read from the lens's polar frame, and the
    elasticity of the count with respect to the radius drawn on the radius
    control, positioned against the local point-pattern statistics it
    re-expresses (\cref{sec:within,sec:elasticity}).
  \item The lens-to-glyphmap continuum as one object (\cref{sec:continuum}).
  \item An open-source implementation (\cref{sec:implementation}) and the
    design of the controlled study the central claim requires
    (\cref{sec:evaluation}).
\end{itemize}

\todo{Teaser figure: category ring $\to$ necklace $\to$ unrolled with leaders
$\to$ field, from \texttt{figures/}.}
